%%%%%%%%%%%%%%%%%
% Lösung von Nazeri et al. 
%%%%%%%%%%%%%%%%%
\section{CPU Results}
\lstset{%
    language         = C++,
    inputencoding    = utf8,
    extendedchars=true,
    literate=%
    {ő}{{{\H{o}}}}1
        {Ő}{{{\H{O}}}}1
        {ö}{{\"o}}1 {ä}{{\"a}}1 {ü}{{\"u}}1 
        {Ö}{{\"O}}1 {Ä}{{\"A}}1 {Ü}{{\"U}}1 
        {ß}{{\ss}}1 
         {ű}{{\H{u}}}1 {Ű}{{\H{U}}}1 % Hungarian fix
        {á}{{\'a}}1 {é}{{\'e}}1 {í}{{\'i}}1 {ó}{{\'o}}1 {ú}{{\'u}}1,
    basicstyle       = \scriptsize\ttfamily,
    keywordstyle     = \bfseries\color{blue},
    stringstyle      = \color{magenta},
    commentstyle     = \color{ForestGreen},
    showstringspaces = false,
}
\lstset{ 
  basicstyle=\scriptsize\ttfamily\linespread{0.7}\selectfont,        % the size of the fonts that are used for the code
  breakatwhitespace=false,         % sets if automatic breaks should only happen at whitespace
  breaklines=true,                 % sets automatic line breaking
  captionpos=b,                    % sets the caption-position to bottom
  %commentstyle=\itshape\color{red},    % comment style
 % deletekeywords={...},            % if you want to delete keywords from the given language
%  escapeinside={\%*}{*)},          % if you want to add LaTeX within your code
  extendedchars=true,              % lets you use non-ASCII characters; for 8-bits encodings only, does not work with UTF-8
  %firstnumber=1000,                % start line enumeration with line 1000
  frame=single,	                   % adds a frame around the code
  keepspaces=true,                 % keeps spaces in text, useful for keeping indentation of code (possibly needs columns=flexible)
 % keywordstyle=\color{blue},       % keyword style
  %language=Octave,                 % the language of the code
  %morekeywords={*,...},            % if you want to add more keywords to the set
  numbers=none,                    % where to put the line-numbers; possible values are (none, left, right)
  numbersep=5pt,                   % how far the line-numbers are from the code
  %numberstyle=\tiny\color{mygray}, % the style that is used for the line-numbers
  rulecolor=\color{black},         % if not set, the frame-color may be changed on line-breaks within not-black text (e.g. comments (green here))
  showspaces=false,                % show spaces everywhere adding particular underscores; it overrides 'showstringspaces'
  showstringspaces=false,          % underline spaces within strings only
  showtabs=false,                  % show tabs within strings adding particular underscores
  stepnumber=5,                    % the step between two line-numbers. If it's 1, each line will be numbered
 % stringstyle=\color{mymauve},     % string literal style
  tabsize=1 ,	                   % sets default tabsize to 2 spaces
  title=\lstname                   % show the filename of files included with \lstinputlisting; also try caption instead of title
%  caption={~/Downloads/Kunoth-21/julia/Blatt07/qrsvd.jl}
}
%%%%%%%%%%%%%%%%%
% \subsection{Evaluierung}
%%%%%%%%%%%%%%%%%
\subsection{Evaluierung}
Ein Knoten $v_i \in V$ gilt als verbrannt, wenn er zu irgendeinem Zeitpunkt des Brennprozesses entzündet wurde. 
Entweder wird er direkt angezündet oder er fängt Feuer durch einen direkten Nachbarn, der bereits brennt. 
Am Ende des Prozesses wird der Abstand von Knoten $v_i$ zum nächsten verbrannten Knoten als Restabstand $d_i$ bezeichnet. 
Für jeden verbrannten Knoten gilt $d_i = 0$. 
Eine valide Brennsequenz zeichnet sich dadurch aus, dass alle Knoten $v_i \in V$ vollständig verbrannt werden. 
Der Preis einer Brennsequenz $S = (v_1, v_2, \ldots, v_b)$ wird als die Summe der quadrierten Restabstände aller Knoten definiert \cite[6]{NAZERI2023}: 
\begin{align}
\label{eq:Cost}
p(S) = \sum_{v_i \in V} d_i^2,
\end{align}
wobei $d_i$ wie folgt berechnet werden kann:   
\begin{align}
\label{eq:d-i}
d_i = 
\begin{cases} 
  0 & \text{falls Knoten } v_i \text{ verbrannt ist,} \\
  \min_{1 \leq j \leq b} \, \{\text{dist}(v_i, v_j) - (b - j) \} & \text{sonst.}
\end{cases}
\end{align} 

Aufgrund der spezifischen Eigenschaften des GBP werden dafür neue Crossover- und Mutationsoperatoren eingeführt, die auf folgenden Eigenschaften basieren \cite[7]{NAZERI2023}: 
\begin{itemize}
\item Die Reihenfolge der Brandherde wird durch diese Operationen nicht verändert. 
\item Die Nachbarschaften der jeweiligen Brandherde werden durch diese Operationen exploriert.
\item Die Exploration der Nachbarschaft eines Knotens soll von seiner Zentralität abhängen.
\end{itemize}

